% !TeX root = ../../Skript.tex
\cohead{\Large\textbf{Rechenoperationen}}
\fakesubsection{Rechenoperationen}
Vergegenwärtigen wir uns nochmals die auf den reellen Zahlen bekannten Rechenoperationen und ihre Gesetze:

\textbf{Addition}

Für zwei beliebige Zahlen \(a,\ b\in\R\) haben wir die Addition mit dem Rechenzeichen \(+\) definiert. Dabei gilt

\begin{itemize}
	\item Das Ergebnis der Addition ist ebenfalls wieder eine reelle Zahl. \(a+b=c\in\R\).
	\item Die Addition ist kommutativ: \(a+b=b+a\)
	\item Es gibt ein Neutralelement der Addition, das wir 0 nennen für das gilt: \(a+0=a\) für alle \(a\in\R\)
	\item Für alle \(a\in\R\) gibt es ein inverses Element in den reellen Zahlen, das wir \(-a\in\R\) nennen, so dass \(a+(-a)=0\) gilt.
	\item Für die Addition gilt das Assoziativgesetz: \((a+b)+c=a+(b+c)\) für alle \(a,\ b,\ c\in\R\)
\end{itemize}

\textbf{Multiplikation}

Für zwei beliebige Zahlen \(a,\ b\in\R\) haben wir die Multiplikation mit dem Rechenzeichen \(\cdot\) definiert. Dabei gilt

\begin{itemize}
	\item Das Ergebnis der Multiplikation ist ebenfalls wieder eine reelle Zahl. \(a\cdot b=c\in\R\).
	\item Die Multiplikation ist kommutativ: \(a\cdot b=b\cdot a\)
	\item Es gibt ein Neutralelement der Multiplikation, das wir 1 nennen für das gilt: \(1\cdot a=a\) für alle \(a\in\R\)
	\item Für alle \(a\in\R\setminus{0}\) gibt es ein inverses Element in den reellen Zahlen, das wir \(a^{-1}=\frac{1}{a}\in\R\) nennen, so dass \(a\cdot a^{-1}=1\) gilt.
	\item Für die Multiplikation gilt das Assoziativgesetz: \((a\cdot b)\cdot c=a\cdot (b\cdot c)\) für alle \(a,\ b,\ c\in\R\)
\end{itemize}

Für die Verknüpfung der beiden Operationen gilt das Distributivgesetz:

\(a\cdot (b+c)=a\cdot b+a\cdot c\)